% Hodina 68, rozpis řádek 68. Zdroje: učebnice
% ucebnice/02_Funkce_a_rovnice/Linearni_funkce_a_rovnice/Lineární funkce.tex,
% oddíly grafy lineárních funkcí a průsečíky; Realisticky.cz lokální
% 020110_Linearni_funkce_III/020110_Lekce.pdf – metodika porovnání více
% přímek v jedné souřadnicové soustavě.
% Úlohy a graf vlastní tvorba. Návaznost na hodinu 69: pololetní písemná práce.
\documentclass[11pt,a4paper]{article}
% Spolecny styl pro ucitelske pripravy MAT (Work/MAT)
% Verze 1.1, 2026-10-23
\NeedsTeXFormat{LaTeX2e}
\ProvidesPackage{priprava}[2026/10/23 v1.1 Sablona ucitelskych priprav MAT]

% Zakladni typografie a cestina
\RequirePackage[T1]{fontenc}
\RequirePackage[utf8]{inputenc}
\RequirePackage[czech]{babel}
\RequirePackage{lmodern}
\RequirePackage{microtype}
\RequirePackage{amsmath,amssymb}
\RequirePackage{xcolor}
\RequirePackage{enumitem}
\RequirePackage[most]{tcolorbox}
\RequirePackage{tabularx,booktabs,array}
\RequirePackage{fancyhdr}
\RequirePackage{needspace}
\RequirePackage{geometry}
\RequirePackage[hidelinks]{hyperref}

\geometry{a4paper,top=17mm,bottom=18mm,left=23mm,right=23mm}
\definecolor{notegray}{gray}{0.39}
\setlength{\parindent}{0pt}
\setlength{\parskip}{2.5pt plus 1pt minus 0.5pt}
\setlist[itemize]{topsep=2pt,itemsep=1pt,parsep=0pt,leftmargin=1.45em}
\setlist[enumerate]{topsep=2pt,itemsep=2pt,parsep=0pt,leftmargin=1.65em}
\setcounter{secnumdepth}{0}
\renewcommand{\arraystretch}{1.12}

% Zahlavi: vlevo rocnik | globalni tema, vpravo cislo: nazev hodiny
% Zapati zustava zcela prazdne.
\pagestyle{fancy}
\fancyhf{}
\renewcommand{\headrulewidth}{0.35pt}
\renewcommand{\footrulewidth}{0pt}
\setlength{\headheight}{14pt}
\newcommand{\nastavPripravu}[4]{%
  \fancyhead[L]{\small #1\enspace\textbar\enspace #2}%
  \fancyhead[R]{\small #3: #4}%
}

% Minimalisticke cernobile boxy. Kazda definice zvlast.
\tcbset{commonlesson/.style={enhanced,breakable=false,colback=white,colframe=black,
  colbacktitle=white,coltitle=black,boxrule=0.45pt,
  arc=0.8mm,outer arc=0.8mm,left=3mm,right=3mm,top=1.6mm,bottom=1.6mm,
  toptitle=1mm,bottomtitle=0.7mm,fonttitle=\bfseries,
  attach boxed title to top left={xshift=2.5mm,yshift=-1.15mm},
  boxed title style={colback=white,colframe=white,boxrule=0pt,left=1mm,right=1mm}}}
\newtcolorbox{zadaniBox}[1][]{commonlesson,title={Zadání},#1}
\newtcolorbox{prikladBox}[1][]{commonlesson,title={Příklad},#1}
\newtcolorbox{definiceBox}[1][]{commonlesson,title={Definice},#1}

% Spolecny box pro zadani a jeho reseni, bez mezery mezi castmi.
% Pouziti:
%   \begin{prikladReseniBox}
%     Zadani ulohy.
%     \tcblower
%     \begin{align*} ... \end{align*}
%   \end{prikladReseniBox}
% Spodni cast ma jemny sedy podklad a delici linku; lze v ni sazet
% vice radku matematiky. Obe casti maji jen jeden spolecny ramecek.
\tcbset{solutionlesson/.style={
  commonlesson,
  enhanced jigsaw,
  breakable,
  colbacklower=black!3,
  segmentation style={solid,draw=black!25,line width=0.35pt},
  fontlower=\small\color{notegray},
  before lower={\noindent\textbf{Řešení:}\enspace},
  leftlower=3mm,rightlower=3mm
}}
\newtcolorbox{prikladReseniBox}[1][]{solutionlesson,title={Příklad},#1}
\newtcolorbox{zadaniReseniBox}[1][]{solutionlesson,title={Zadání},#1}

% Pomucky pro autora pripravy: sedive poznamky; \reseni pro starsi pripravy
\newcommand{\poznamka}[1]{\par\begingroup\small\color{notegray}\itshape
  \textnormal{\textbf{Poznámka pro učitele.}}\enspace #1\par\endgroup}
\newcommand{\reseni}[1]{\par\begingroup\small\color{notegray}
  \textbf{Očekávané řešení:}\enspace #1\par\endgroup}
\newcommand{\blok}[1]{\needspace{5\baselineskip}\subsection*{#1}}
\newcommand{\vyrok}[1]{\uv{#1}}

\usepackage{tikz}
\nastavPripravu{1. ročník}{Funkce}{68}{Grafické porovnávání lineárních funkcí}
\begin{document}
\blok{Motivace a návaznost}
Dvě lineární funkce můžeme nakreslit do stejné soustavy a pro každý argument porovnávat jejich výšky. Průsečík grafů je situace, kdy obě funkce při stejném vstupu dávají stejný výstup. Předchozí hodina ukázala praktické významy parametrů; nyní je spojíme s grafickým a algebraickým porovnáním.
\begin{definiceBox}
\textbf{Průsečík grafů} funkcí $f$ a $g$ je bod $(x;y)$, pro který $x\in D(f)\cap D(g)$ a $y=f(x)=g(x)$. Jeho argument získáme řešením rovnice $f(x)=g(x)$; výsledek musí patřit do společného definičního oboru.
\end{definiceBox}
\blok{Dvě různě skloněné přímky}
\begin{prikladReseniBox}
Porovnejte $f(x)=x+1$ a $g(x)=-x+5$ na $\mathbb R$. Určete jejich průsečík a rozhodněte, která má vyšší hodnotu vlevo a vpravo od něj.
\tcblower
Rovnost hodnot dává $x+1=-x+5$, tedy $2x=4$ a $x=2$. Obě funkce tam mají hodnotu $3$, takže průsečík je $(2;3)$. Rozdíl $f(x)-g(x)=2x-4=2(x-2)$ je pro $x<2$ záporný a pro $x>2$ kladný. Vlevo je výše graf $g$, vpravo graf $f$. Kontrola dosazením: při $x=0$ jsou hodnoty $1$ a $5$; při $x=3$ jsou $4$ a $2$.
\end{prikladReseniBox}
\begin{center}
\begin{tikzpicture}[x=.8cm,y=.6cm,>=stealth]
\draw[step=1,gray!25,very thin] (-1,-1) grid (5,6);
\draw[->] (-1.3,0)--(5.4,0) node[right] {$x$};\draw[->] (0,-1.3)--(0,6.4) node[above] {$y$};
\foreach \x in {1,2,3,4,5}{\draw (\x,.08)--(\x,-.08) node[below] {\small $\x$};}
\foreach \y in {1,3,5}{\draw (.08,\y)--(-.08,\y) node[left] {\small $\y$};}
\draw[very thick] (-1,0)--(5,6) node[above] {\small $f$};
\draw[very thick,dashed] (-1,6)--(5,0);
\node[above right=3pt] at (4,1) {\small $g$};
\fill (2,3) circle (2.5pt);\node[right=8pt] at (2,3) {\small $(2;3)$};
\end{tikzpicture}
\end{center}
Plná a čárkovaná čára zůstávají rozlišitelné i při černobílém tisku. Na svislici $x=2$ se setkají ve stejné výšce; jinde lze vyšší hodnotu odečítat přibližně z grafu a ověřit výpočtem. Při nejistotě se znaménkem rozdílu $f(x)-g(x)$ nejprve porovnejte výšky obou grafů v konkrétním bodě na každé straně průsečíku.
\blok{Rovnoběžné a totožné grafy}
\begin{prikladReseniBox}
Porovnejte $p(x)=2x+1$ a $q(x)=2x-3$ na $\mathbb R$. Mají průsečík?
\tcblower
Směrnice jsou stejné, grafy jsou rovnoběžné. Rovnost $2x+1=2x-3$ by vyžadovala $1=-3$, což není možné. Průsečík nemají. Pro každé $x$ platí $p(x)-q(x)=4$, takže graf $p$ leží stále o $4$ výše.
\end{prikladReseniBox}
\begin{prikladReseniBox}
Porovnejte $u(x)=3(x-1)+2$ a $v(x)=3x-1$ na $\mathbb R$.
\tcblower
Po roznásobení $u(x)=3x-3+2=3x-1=v(x)$ pro všechna reálná $x$. Grafy splývají a každý jejich bod je společný. Kontrola například při $x=0$: obě hodnoty jsou $-1$.
\end{prikladReseniBox}
\poznamka{Rozliš tři možnosti: různá směrnice znamená jeden průsečík na $\mathbb R$; stejná směrnice s různým úsekem žádný; shodné oba koeficienty totožné grafy. U omezených oborů musí vypočtený argument průsečíku ležet v obou oborech.}
\blok{Samostatné procvičení}
\begin{zadaniReseniBox}
\begin{enumerate}[label=\alph*)]
\item Sestrojte do jedné soustavy grafy $f(x)=2x-1$ a $g(x)=x+2$ (např. z bodů pro $x=0$ a $x=2$). Z grafu odečtěte průsečík i intervaly, kde je která funkční hodnota vyšší. Výsledek ověřte algebraicky.
\item Mají průsečík $h(x)=-x+3$ a $k(x)=-x-2$?
\item Funkce $m(x)=x+1$ je dána jen pro $x\in\langle0;1\rangle$, funkce $n(x)=-x+5$ pro $x\in\langle0;3\rangle$. Protínají se jejich skutečné grafy?
\end{enumerate}
\tcblower
\begin{enumerate}[label=\alph*)]
\item Z předpisů sestrojíme $f$ přes $(0;-1),(2;3)$ a $g$ přes $(0;2),(2;4)$. V grafu čteme průsečík $(3;5)$; pro $x<3$ leží výše graf $g$, pro $x>3$ graf $f$, v $x=3$ se hodnoty rovnají. Algebraicky $2x-1=x+2$ dává $x=3$ a obě hodnoty $5$. Rozdíl $f(x)-g(x)=x-3$ potvrzuje pořadí pro všechny argumenty. Kontrola: při $x=0$ jsou hodnoty $-1,2$ a při $x=4$ hodnoty $7,6$.
\item Ne: směrnice jsou obě $-1$, ale úseky $3$ a $-2$ se liší. Rozdíl hodnot je stále $5$.
\item Rovnice předpisů dává $x=2$, ale $2\notin D(m)=\langle0;1\rangle$. Skutečné grafy nemají společný bod, i když prodloužené přímky se protínají v $(2;3)$.
\end{enumerate}
\end{zadaniReseniBox}
\blok{Volitelné rozšíření a rezerva}
\begin{zadaniReseniBox}
Dva tarify stojí $A(t)=100+4t$ a $B(t)=40+7t$ Kč pro $t\ge0$ minut. Při jaké době jsou ceny stejné a který tarif je levnější po $10$ minutách?
\tcblower
$100+4t=40+7t$ znamená $60=3t$, tedy $t=20$ min; oba stojí $180$ Kč. Po $10$ min je $A(10)=140$ Kč a $B(10)=110$ Kč, levnější je $B$. Před $20$ min leží jeho graf níže, po $20$ min je levnější $A$. Model předpokládá ceny lineární pro uvedený čas.
\end{zadaniReseniBox}
\blok{Domácí práce}
\begin{zadaniReseniBox}
\begin{enumerate}[label=\alph*)]
\item Najděte průsečík $f(x)=-2x+6$ a $g(x)=x$ a ověřte jej dosazením.
\item Porovnejte $p(x)=\frac12x+1$ a $q(x)=\frac12x-2$.
\item Dva náklady jsou $C_1(n)=50n+300$ a $C_2(n)=70n+100$ Kč, $n\in\mathbb N_0$. Pro kolik kusů jsou stejné?
\end{enumerate}
\tcblower
\begin{enumerate}[label=\alph*)]
\item $-2x+6=x$ dává $x=2$, hodnota $2$, průsečík $(2;2)$. Kontrola $f(2)=2=g(2)$.
\item Směrnice jsou stejné, rozdíl $p(x)-q(x)=3$. Přímky jsou rovnoběžné a $p$ leží všude o $3$ výše.
\item $50n+300=70n+100$ dává $200=20n$, tedy $n=10$ kusů. Obě ceny jsou $800$ Kč; $10\in\mathbb N_0$.
\end{enumerate}
\end{zadaniReseniBox}
\poznamka{Jádro 45 min: návaznost 4, dvě přímky a jejich graf 13, rovnoběžnost a totožnost 8, samostatný pokus 9, společná kontrola 8, organizace 3 min. Rezerva 6–8 min, domácí 10–15 min. Následuje hodina 69: pololetní písemná práce z algebry a základů funkcí; zde nepřidávej novou látku mimo grafické porovnávání.}
\end{document}
